\documentclass[11pt,a4paper]{article}


% ===============================
% Chinese + fonts (XeLaTeX/LuaLaTeX)
% ===============================
\usepackage[UTF8,scheme=plain,fontset=none]{ctex} % it handles fontspec loading for Chinese
\setmainfont{TeX Gyre Termes} % like Times New Roman \textrm{...}
\setsansfont{TeX Gyre Heros}  % like Arial \textsf{...}
\setmonofont{TeX Gyre Cursor} % like Consolas \texttt{...}

% CJK fonts (Noto fonts from Google)
\setCJKmainfont{Noto Serif CJK SC}  % Song Ti
\setCJKsansfont{Noto Sans CJK SC}   % Hei Ti \textsf{...}
\setCJKmonofont{Noto Sans Mono CJK SC} % Mono Font \texttt{...}


% ===============================
% page layout
% ===============================
\usepackage[a4paper,margin=2.5cm]{geometry}
\usepackage{microtype}


% ===============================
% math packages
% ===============================
\usepackage{amsmath,amssymb,amsfonts,amsthm}
\usepackage{mathtools}
\usepackage{bm}


% ===============================
% graphics
% ===============================
\usepackage{graphicx}
\graphicspath{{Images/}{./}}
\usepackage{subcaption}
\usepackage{wrapfig} % wrap text around figures (\begin{wrapfigure}...)


% ===============================
% colors
% ===============================
\usepackage{xcolor}  % color text (\textcolor{red}{...})


% ===============================
% drawing packages
% ===============================
\usepackage{tikz}
\usetikzlibrary{
    arrows.meta,     % more modern arrow tips
    calc,            % coordinate calculations
    positioning,     % relative position（right=of ...）
    decorations.pathmorphing, % 
    patterns,        %
    shapes.geometric,%
    fit              % 
}


% ===============================
% tables
% ===============================
\usepackage{booktabs} % better tables
\usepackage{array}    % better arrays and tabulars
\usepackage{multirow} % multirow cells in tables


% ===============================
% bibliography (recommended)
% ===============================
\usepackage[
    backend=biber,
    style=authoryear,
    sorting=nyt
    ]{biblatex} % \cite{key} for citation
                % \printbibliography to print the reference list
\addbibresource{references.bib}


% ===============================
% hyperlinks 
% ===============================
\usepackage{hyperref}
\hypersetup{
    colorlinks=true,
    linkcolor=blue,
    citecolor=blue,
    urlcolor=blue
}


% ===============================
% outside files
% ===============================
% \newcommand ---- define new command
% ===============================
% Basic sets of numbers
% ===============================
\newcommand{\R}{\mathbb{R}}
\newcommand{\N}{\mathbb{N}}
\newcommand{\Z}{\mathbb{Z}}
\newcommand{\Q}{\mathbb{Q}}

% ===============================
% Differential operators and constants
% ===============================
\newcommand{\dd}{\,\mathrm{d}}
\newcommand{\ee}{\mathrm{e}}
\newcommand{\ii}{\mathrm{i}}

% ===============================
% Vector calculus operators
% ===============================
\newcommand{\grad}{\nabla}
\newcommand{\diver}{\nabla\cdot}
\newcommand{\curl}{\nabla\times}


% ===============================
% Function spaces
% ===============================
\newcommand{\Hone}{H^1(\Omega)}
\newcommand{\Ltwo}{L^2(\Omega)}
\newcommand{\Hdiv}{H(\text{div},\Omega)}
\newcommand{\Hcurl}{H(\text{curl},\Omega)}


% ===============================
% Common operators (amsmath required)
% Put this in mathcmd.tex, and ensure amsmath is loaded before % \newcommand ---- define new command
% ===============================
% Basic sets of numbers
% ===============================
\newcommand{\R}{\mathbb{R}}
\newcommand{\N}{\mathbb{N}}
\newcommand{\Z}{\mathbb{Z}}
\newcommand{\Q}{\mathbb{Q}}

% ===============================
% Differential operators and constants
% ===============================
\newcommand{\dd}{\,\mathrm{d}}
\newcommand{\ee}{\mathrm{e}}
\newcommand{\ii}{\mathrm{i}}

% ===============================
% Vector calculus operators
% ===============================
\newcommand{\grad}{\nabla}
\newcommand{\diver}{\nabla\cdot}
\newcommand{\curl}{\nabla\times}


% ===============================
% Function spaces
% ===============================
\newcommand{\Hone}{H^1(\Omega)}
\newcommand{\Ltwo}{L^2(\Omega)}
\newcommand{\Hdiv}{H(\text{div},\Omega)}
\newcommand{\Hcurl}{H(\text{curl},\Omega)}


% ===============================
% Common operators (amsmath required)
% Put this in mathcmd.tex, and ensure amsmath is loaded before \input{mathcmd}
% ===============================
% --- Linear algebra / matrix analysis ---
\DeclareMathOperator{\rank}{rank}
\DeclareMathOperator{\tr}{tr}
\DeclareMathOperator{\diag}{diag}
\DeclareMathOperator{\trace}{trace}


% ===============================
% --- Linear algebra / matrix analysis ---
\DeclareMathOperator{\rank}{rank}
\DeclareMathOperator{\tr}{tr}
\DeclareMathOperator{\diag}{diag}
\DeclareMathOperator{\trace}{trace}



% \newtheorem{<env>}{<Title>} ---- define a new environment
% \newtheorem{<env>}{<Title>}[<reset>]
% \newtheorem{<env>}[<share>]{<Title>}

% ===============================
% Theorem environments
% ===============================
\theoremstyle{plain}

\newtheorem{theorem}{Theorem}[section]
\newtheorem{lemma}[theorem]{Lemma}
\newtheorem{proposition}[theorem]{Proposition}
\newtheorem{corollary}[theorem]{Corollary}

% ===============================
% Definition environments
% ===============================
\theoremstyle{definition}

\newtheorem{definition}[theorem]{Definition}
\newtheorem{assumption}[theorem]{Assumption}
\newtheorem{example}[theorem]{Example}

% ===============================
% Remark environments
% ===============================
\theoremstyle{remark}

\newtheorem{remark}[theorem]{Remark}
\newtheorem{note}[theorem]{Note}




\title{Article Demo: Using \texttt{preamble.tex}, \texttt{mathcmd.tex}, \texttt{theorem.tex}}
\author{song wu}
\date{\today}

% ==============================
\begin{document}
% ==============================
\maketitle

\begin{abstract}
This demo exercises most packages and all commands/environments defined in the provided files:
\texttt{fontspec}, page layout, math tools, figures (\texttt{graphicx/subcaption/wrapfig}),
TikZ libraries, tables, \texttt{hyperref}, and \texttt{biblatex}.
We cite \textcite{lamport1994} and also \parencite{einstein1905}.
\end{abstract}

\tableofcontents
\newpage

\section{Math macros from \texttt{mathcmd.tex}}
\subsection{Basic sets and constants}
这是中文测试。基础数学符号测试。
We use $\R,\N,\Z,\Q$ and $\ee,\ii$:
\[
  \ee^{\ii\pi}+1=0,\qquad
  \int_0^1 x^2\,\dd x=\frac13.
\]

\subsection{Vector calculus operators}
Let $u(x,y)=x^2+y^2$ and $\bm{v}(x,y)=(x,-y)^\top$.
Then we can write
\[
  \grad u,\quad
  \diver \bm{v},\quad
  \curl \bm{v}.
\]
(Here we mainly show the macros/typesetting; \\
the actual values depend on the dimension/convention.)

\subsection{Function spaces}
We reference the function spaces:
\[
  u\in\Hone,\quad
  f\in\Ltwo,\quad
  \bm{u}\in\Hdiv,\quad
  \bm{w}\in\Hcurl.
\]

\subsection{Linear algebra operators}
Declared operators:
\[
  \rank(A),\quad \tr(A),\quad \diag(1,2,3),\quad \trace(A).
\]
Example:
\[
  A=\begin{pmatrix}2&0\\0&1\end{pmatrix}
  \Rightarrow \rank(A)=2,\, \tr(A)=3,\ \trace(A)=3,\,\det(A)=2.
\]

\section{Theorem-like environments from \texttt{theorem.tex}}

\begin{definition}[SPD matrix]
A matrix $A\in\R^{n\times n}$ is symmetric positive definite (SPD) if
$A=A^\top$ and $\bm{x}^\top A\bm{x}>0$ for all $\bm{x}\neq 0$.
\end{definition}

\begin{assumption}
Assume $A$ is SPD.
\end{assumption}

\begin{lemma}
If $A$ is SPD, then all eigenvalues of $A$ are positive.
\end{lemma}

\begin{theorem}
If $A$ is SPD, then $\tr(A)>0$.
\end{theorem}

\begin{proof}
SPD implies eigenvalues $\lambda_i>0$. Then $\tr(A)=\sum_i \lambda_i>0$.
For general LaTeX conventions, see \textcite{lamport1994}.
\end{proof}

\begin{proposition}
If $A$ is SPD and $D=\diag(d_1,\dots,d_n)$ with all $d_i>0$, then $DAD$ is SPD.
\end{proposition}

\begin{corollary}
If $A$ is SPD, then $\trace(A^{-1})>0$.
\end{corollary}

\begin{example}
Let $A=\diag(2,1)$. Then $A$ is SPD and $\tr(A)=3$.
\end{example}

\begin{remark}
All these environments share the same counter (reset by section), \\
as configured in \texttt{theorem.tex}.
\end{remark}

\begin{note}
This document uses \texttt{fontspec}, so compile with XeLaTeX or LuaLaTeX.
\end{note}

\section{Graphics: \texttt{wrapfig}, \texttt{subcaption}, \texttt{graphicx}}

\subsection{Wrapfigure demo}
\begin{wrapfigure}{r}{0.35\textwidth}
    \centering
    \begin{tikzpicture}[scale=0.95]
        \draw[->] (-0.2,0) -- (2.6,0) node[right] {$x$};
        \draw[->] (0,-0.2) -- (0,2.2) node[above] {$y$};
        \draw[thick,blue,domain=0:2.3,samples=60] plot (\x,{0.4*\x*\x});
        \node[draw,rounded corners,fill=blue!10] at (1.6,1.7) {$y=0.4x^2$};
    \end{tikzpicture}
    \caption{A wrapped TikZ figure.}
\end{wrapfigure}

This paragraph demonstrates \texttt{wrapfig}. 
Add a few lines of text so the wrapping is clearly visible.
Here is some more text to illustrate how the figure wraps around the text.
The figure is placed on the right side of the page, and the text flows around it.
You can adjust the width of the figure and the placement (left or right) as needed.
Continue adding text to see how it interacts with the wrapped figure.
This is useful for including small figures or diagrams without interrupting the flow of the text.

\subsection{Subcaption demo (four figures)}

\begin{figure}[htbp] % 2 x 2
    \centering
    \begin{subfigure}[b]{0.48\textwidth}
        \centering
        \includegraphics[width=\linewidth]{call-of-duty.jpeg}
        \caption{call of duty}
        \label{fig:a}
    \end{subfigure}\hfill
    \begin{subfigure}[b]{0.48\textwidth}
        \centering
        \includegraphics[width=\linewidth]{f1-the-movie.jpg}
        \caption{f1 the movie}
        \label{fig:b}
    \end{subfigure}

    \medskip  % 行间距（推荐）

    \begin{subfigure}[b]{0.48\textwidth}
        \centering
        \includegraphics[width=\linewidth]{naraka-bladepoint.jpeg}
        \caption{naraka bladepoint}
        \label{fig:c}
    \end{subfigure}\hfill
    \begin{subfigure}[b]{0.48\textwidth}
        \centering
        \includegraphics[width=\linewidth]{your-name-shooting.jpg}
        \caption{your name shooting}
        \label{fig:d}
    \end{subfigure}

    \caption{A 2×2 layout of four images}
    \label{fig:2x2}
\end{figure}




\subsection{TikZ libraries demo: calc, patterns, shapes, decorations}

\begin{figure}[htbp]
    \centering
    \begin{subfigure}[b]{0.47\textwidth}
        \centering
        \begin{tikzpicture}
            \draw[thick] (0,0) rectangle (4,2);
            \draw[very thick,-{Stealth[length=3mm]},red] (0.6,0.6) -- (3.4,1.4) node[right] {$\bm{u}$};
            \node[below] at (2,0) {box};
        \end{tikzpicture}
        \caption{Arrow tip from \texttt{arrows.meta}.}
    \end{subfigure}\hfill
    \begin{subfigure}[b]{0.47\textwidth}
        \centering
        \begin{tikzpicture}
            \node[draw,circle,minimum size=16mm,fill=green!10] (A) at (0,0) {$\R^2$};
            \node[draw,rounded corners,minimum width=30mm,minimum height=12mm,
                right=18mm of A,fill=yellow!10] (B) {$\Hcurl$};
            \draw[-{Stealth[length=3mm]},thick] (A) -- node[above] {map} (B);
            % fit library demo
            \node[draw,dashed,fit=(A)(B),inner sep=4mm,label=above:{fit}] {};
        \end{tikzpicture}
        \caption{Positioning + fit.}
    \end{subfigure}
    \caption{Two subfigures without external image files.}
\end{figure}


\section{Colors: \texttt{xcolor}}
We can also show colored text:\textcolor{red}{this is red}, \textcolor{blue}{this is blue}.



\section{Tables: \texttt{booktabs}, \texttt{array}, \texttt{multirow}}

\begin{table}[htbp]
    \centering
    \renewcommand{\arraystretch}{1.15}
    \begin{tabular}{>{\bfseries}l c c}
        \toprule
        Item & Example & Note \\
        \midrule
        \multirow{2}{*}{Trace} & $\tr(A)$ & operator \\
                               & $\trace(A)$ & alias \\
        Rank & $\rank(A)$ & operator \\
        Space & $\Hdiv,\ \Hcurl$ & macros \\
        \bottomrule
    \end{tabular}
    \caption{A table demonstrating booktabs + array + multirow.}
\end{table}

\section{Hyperlinks and bibliography}

A URL: \url{https://www.latex-project.org}. A clickable link: \href{https://ctan.org}{CTAN}.

We cite \parencite{trefethen1997} for numerical linear algebra.

\printbibliography

% ==============================
\end{document}
% ==============================
